\section{Performance and memory}
\label{sec:perf}

\subsection{Cost model}

The per-chunk cost is
\begin{equation}
  C_{\text{chunk}} =
  C_{\text{flatten}} + C_{\text{upload}} + C_{\text{walk}} +
  C_{\text{read}} + C_{\text{reserve}} + C_{\text{finalize}} +
  C_{\text{publish}}.
\end{equation}
The flatten and upload scales with the number of nodes in the
neighbourhood; the walk scales with the number of scan cells and active
edges; the finalize scales with the index and vertex counts.  The host
reads only six integers between the two submissions.

\subsection{Measured topology rates}

The debug profiler in \code{OctreeTopologyCompute} reports cumulative
averages every 256 chunks when \code{VULKAN_SETUP_DIAG=1}.  On the reference
scene the observed windows were:

\begin{center}
\begin{tabular}{@{}rrrr@{}}
\toprule
Chunks & avg ms/chunk & avg cells & avg nodes\\
\midrule
256  & 8.3  & 1031 & 3761\\
512  & 9.3  & 911  & 3440\\
768  & 10.0 & 885  & 3382\\
1280 & 18.9 & 788  & 3053\\
\bottomrule
\end{tabular}
\end{center}

The later window is slower despite fewer cells and nodes because the
workers contend with the render thread and with descriptor/buffer growth.
The uncontended floor reported for the last window of an earlier capture
was about 4.2 ms/chunk.  The difference is a convoy effect, not a
super-linear algorithm.

\begin{figure}[htbp]
\centering
\begin{tikzpicture}
\begin{axis}[
  width=0.78\linewidth,height=5.2cm,
  xlabel={chunk window (cumulative)},ylabel={ms/chunk},
  xtick={256,512,768,1280},
  ymin=0,ymax=22,
  grid=major,grid style={draw=ink!12},
  legend style={font=\scriptsize,at={(0.98,0.97)},anchor=north east},
  tick label style={font=\scriptsize},
  label style={font=\small}
]
\addplot[mark=*,thick,accent] coordinates {(256,8.3) (512,9.3) (768,10.0) (1280,18.9)};
\addplot[mark=square*,thick,accent2,dashed] coordinates {(256,4.2) (512,4.2) (768,4.2) (1280,4.2)};
\legend{observed window, uncontended floor}
\end{axis}
\end{tikzpicture}
\caption{Observed topology cost per chunk.  The floor is the cost when the
worker pool is not contending with the render thread or with buffer
growth.}
\label{fig:perf}
\end{figure}

\subsection{Throughput}

The walk is one invocation per scan cell.  For a chunk with about 900 scan
cells and 3.4k nodes, the measured 8--19 ms includes the host-side flatten
and upload, the two submissions, and the fence wait.  The GPU dispatch
itself is a small fraction: the adaptive walk is branch-heavy and the
per-invocation arrays (break list, stack, memo) are larger than the scan
loop.  The cost is dominated by the branch divergence of the edge
classification and by the break collection, not by the QEF solve.

\subsection{Memory budget}

The reference configuration reserves:
\begin{center}
\begin{tabular}{@{}lrr@{}}
\toprule
Resource & Size & Notes\\
\midrule
Solid vertex pool & 768 MB & 64 B $\times$ 12M elements\\
Solid index pool & 192 MB & 4 B $\times$ 48M elements\\
Water vertex pool & 96 MB & 64 B $\times$ 1.5M elements\\
Water index pool & 96 MB & 4 B $\times$ 24M elements\\
Vegetation concat A/B & 256 MB $\times$ 2 & 8M instances $\times$ 32 B\\
Vegetation baked A/B & 32 MB $\times$ 2 & one float per instance\\
Vegetation scratch & 80 MB & 2 $\times$ 10 $\times$ 131k $\times$ 32 B\\
Texture arrays & 720 / 747 MB & five maps, 28 layers\\
\bottomrule
\end{tabular}
\end{center}

The pools are preallocated once.  A publish only carves or frees a span;
it never allocates a Vulkan buffer.  This is why the steady-state path has
zero \code{vmaCreateBuffer} calls and why a streaming burst cannot
fragment the allocator.

\subsection{Streaming behaviour}

Chunks are drained at a controlled rate (16 per frame in the scene
publisher, 10 per frame for vegetation).  A full drain would submit
hundreds of topology builds at once; on a 4 GB iGPU that trips the amdgpu
watchdog.  The rate limit keeps the submit queue bounded and spreads the
GPU work over frames.  The vegetation and fire effects lag the solid
publish by one or two frames; that is acceptable because the chunk is not
visible until its draw entry is published.

\subsection{Failure costs}

The expensive failure is a key-capacity retry: it re-uploads the node
table and re-runs the walk.  The hint (25 keys per scan cell plus 64) was
calibrated to make retries rare.  The hard cap (8M keys, 32 MB) is chosen
so a pathological chunk fails explicitly rather than allocating
unbounded memory.  A pool exhaustion is a sizing bug: the pools are
preallocated for the worst measured scene, and a new scene must bump the
budget deliberately.

\subsection{Interpretation}

The pipeline is not ``free''.  It pays:
\begin{itemize}
\item a per-chunk flatten and upload of the node neighbourhood;
\item two submissions and one fence wait;
\item a small host read of six counters;
\item a per-chunk span reservation and descriptor update.
\end{itemize}
It saves:
\begin{itemize}
\item a GPU$\rightarrow$CPU vertex readback;
\item a CPU-side QEF solve for every vertex;
\item a host-side triangle append and a staging upload;
\item a CPU-to-GPU copy of the index stream.
\end{itemize}
On a discrete GPU the readback and staging are the dominant costs; on an
integrated GPU they are even more expensive because the memory is shared.
The measured topology cost is the price of the split, and it is paid on
worker threads, not on the render thread.
